\begin{explanationbox}
	\textbf{\textcolor{CExplanation}{一句话直觉}}
	点到直线的距离，可以看作过点作直线的垂线，垂线段的长度。

	\medskip
	\textbf{\textcolor{CExplanation}{核心拆解}}
	\begin{description}[style=nextline, leftmargin=2.8em, labelsep=0.5em, itemsep=0.45em]
		\item[\textbf{要点一（分子几何意义）：}]
			分子 $|Ax_0+By_0+C|$ 表示将点坐标代入直线方程后得到的绝对值；当点在直线上时该值为 $0$，点离直线越远绝对值越大。
		\item[\textbf{要点二（分母归一化因子）：}]
			分母 $\sqrt{A^2+B^2}$ 是直线法向量 $(A,B)$ 的模长，起“归一化”作用，保证距离与直线系数尺度无关。
	\end{description}

	\medskip
	\textbf{\textcolor{CExplanation}{几何本质（法向量投影）}}
	点 $P$ 到直线的距离等于向量 $\overrightarrow{PP_0}$（$P_0$ 为直线上任意一点）在直线法向量方向上的投影长度。公式正好是这个投影的绝对值。

	\medskip
	\textbf{\textcolor{CExplanation}{代数意义（条件约束）}}
	公式提供了一个统一的方法，不需要将直线化为斜截式，只需一般式系数即可计算距离，体现了解析几何中代数化的思想。

	\medskip
	\textbf{\textcolor{CExplanation}{考点价值（计算基础）}}
	\begin{itemize}[leftmargin=*, itemsep=0.5em, topsep=0.4em]
		\item \textbf{考法一（直接求距离）：} 给出直线与点，直接代入公式求距离。
		\item \textbf{考法二（求参数或范围）：} 已知距离反求直线系数或点坐标，建立含绝对值的方程。
	\end{itemize}

	\medskip
	\textbf{\textcolor{CExplanation}{顿悟点}}
	\textbf{距离公式的分母 $\sqrt{A^2+B^2}$ 与法向量长度挂钩，分子体现点的“偏离程度”。}

	\medskip
	\textbf{\textcolor{CExplanation}{使用场景}}
	\begin{itemize}[leftmargin=*, itemsep=0.5em, topsep=0.4em]
		\item \textbf{场景一（点线距计算）：} 任意给定直线一般式和点坐标，求点到直线的距离。
		\item \textbf{场景二（平行线距转化）：} 求两平行直线间距离时，可转化为一条直线上的一点到另一条直线的距离。
	\end{itemize}

\end{explanationbox}
